This work set out to rescue SJCET's handwritten weather register and to find out whether a data-driven model could forecast campus weather from it. Both parts have been achieved for the July 2010 -- December 2016 record.

The digitization produced a traceable, quality-flagged daily dataset of 1,549 observed days, with every value linked to its page and 59 of 78 months confirmed against the observers' own rainfall totals. It also documented the register's quirks (swapped columns, blank months, a rarely reset maximum thermometer) and separated our transcription errors from the observers' arithmetic slips. The OCR experiment showed that off-the-shelf handwriting recognition is not yet good enough for this register.

The forecasting experiment answered its central question. A graph neural network pretrained on regional reanalysis and fine-tuned on SJCET forecasts next-day maximum and minimum temperature with 96.6\% and 96.2\% accuracy, with about a third less error than the same network trained on SJCET alone. It beats XGBoost on maximum temperature at every lead day, and it overtakes persistence from day 4. With only 10\% of the local data it is the best of all methods for maximum temperature, but worse than plain climatology for minimum temperature. That negative result points at the next improvement. The station stage helped only when local data was scarcest, and rainfall remains essentially unpredictable at a single gauge.

Three practical lessons are worth recording. First, \textbf{the checksum is worth more than the OCR}: the observers' monthly totals did more for data quality than any recognition model. Second, \textbf{the boring baselines matter}: without persistence and climatology the GNN results would look better than they are. Third, \textbf{the files-as-interface design paid for itself}: developing the model on a stand-in station during transcription meant only one stage had to be rerun once the data arrived.

A cupboard of handwritten registers is not an archive until someone reads it. Read carefully and joined to the abundant regional data that already exists, a small local record becomes both a piece of the campus's climate history and the most valuable ingredient of its own forecast.
